\VerdierPhysicalPageAnchor{0076}%
\VerdierSemanticAnchor{ch1-2-4-2-countable-products-sums}
Supposons maintenant que les produits (resp. sommes directes) dénombrables
soient représentables dans $C$. On pose alors :

\VerdierSemanticAnchor{ch1-2-4-2-eq6}
\begin{equation}
  \begin{aligned}
    \operatorname{comp}_{\Pi}F
      &=\int_n^{\Pi}F^{[n]},\\
    \text{(resp. }\operatorname{comp}_{\Sigma}F
      &=\int_n^{\Sigma}F^{[n]}\text{)}.
  \end{aligned}
  \tag{2.4.2.6}
\end{equation}

\VerdierSemanticAnchor{ch1-2-4-2-pi-sigma-extension}
Le foncteur
$\operatorname{comp}_{\Pi}F$ (resp. $\operatorname{comp}_{\Sigma}F$) :
$\prod_{i\in[n]}\operatorname{comp}(C_i)\longrightarrow
\operatorname{comp}(C)$ est appelé la $\Pi$-extension (resp. la
$\Sigma$-extension) du foncteur $F$ aux complexes.

\VerdierSemanticAnchor{ch1-2-4-2-koszul-twist}
Le foncteur $F$ étant un $\mathbb Z[n]$-foncteur strict et le foncteur
$\int_n^{\Pi}$ (resp. $\int_n^{\Sigma}$) étant un $\mathbb Z[n]$-foncteur
tordu par le cocycle de Koszul relatif à la base ordonnée
$e_1,\ldots,e_n$ de $\mathbb Z[n]$, le foncteur
$\operatorname{comp}_{\Pi}F$ (resp. $\operatorname{comp}_{\Sigma}F$) est
un $\mathbb Z[n]$-foncteur tordu par le cocycle de Koszul relatif à la base
ordonnée $e_1,\ldots,e_n$ de $\mathbb Z[n]$, conformément à la règle des
signes I (1.6.6). On a donc pour tout $i\in[n]$ des isomorphismes
canoniques :

\VerdierSemanticAnchor{ch1-2-4-2-eq7}
\begin{equation}
  \begin{aligned}
    \operatorname{comp}_{\Pi}F(X_1,\ldots,X_i[1],\ldots,X_n)
      &\xrightarrow{\sim}
      \operatorname{comp}_{\Pi}F(X_1,\ldots,X_i,\ldots,X_n)[1],\\
    \text{(resp. }\operatorname{comp}_{\Sigma}F(X_1,\ldots,X_i[1],\ldots,X_n)
      &\xrightarrow{\sim}
      \operatorname{comp}_{\Sigma}F(X_1,\ldots,X_i,\ldots,X_n)[1]\text{)}
  \end{aligned}
  \tag{2.4.2.7}
\end{equation}

\VerdierSemanticAnchor{ch1-2-4-2-canonical-pi}
qui seront notés $p_i$.

\VerdierSemanticAnchor{ch1-2-4-2-anticommutativity}
Ces isomorphismes possèdent les propriétés d'anticommutativité décrites en
(1.6.6.3).

\VerdierSemanticAnchor{ch1-2-4-2-summable-restriction-equality}
Les restrictions des foncteurs $\operatorname{comp}_{\Pi}F$ et
$\operatorname{comp}_{\Sigma}F$ aux objets sommables de
$\prod_{i\in[n]}\operatorname{comp}(C_i)$ sont égales, et cette restriction
peut toujours se définir, sans hypothèses sur la catégorie $C$.

\VerdierSemanticAnchor{ch1-2-4-3}
\subsection*{2.4.3.}
Soient maintenant deux foncteurs multiadditifs
$F,F'\colon\prod_{i\in[n]}C_i\longrightarrow C$ et
$h\colon F\longrightarrow F'$ un morphisme de foncteurs. On en déduit un
morphisme de foncteurs :

\VerdierSemanticAnchor{ch1-2-4-3-induced-morphism}
\begin{equation*}
  h^{[n]}\colon F^{[n]}\longrightarrow F'^{[n]},
\end{equation*}
en posant :

\VerdierSemanticAnchor{ch1-2-4-3-eq1}
\begin{equation}
  \begin{aligned}
    h^{[n]}(X_1,\ldots,X_n)^{\sigma}
      &=h(X_1^{l_1},\ldots,X_n^{l_n}),\\
    \sigma&=\sum_{i\in[n]}l_i e_i\in\mathbb Z[n],\\
    (X_1,\ldots,X_n)&\in
      \operatorname{Ob}\!\left(\prod_{i\in[n]}\operatorname{comp}(C_i)\right).
  \end{aligned}
  \tag{2.4.3.1}
\end{equation}
